\chaptertitle{5}{常见的宽基指数有哪些}

我们已经知道了，指数基金有很多好处，可是市面上的指数基金这么多，不可能全部都投资了，那哪些品种适合长期定投呢？

在接下来的几章，我们就来了解一下。

\subsection{指数的分类方式}
常见指数有哪些呢？指数有两种主要的分类方式。

\subsubsection{宽基指数和行业指数}

有的指数在挑选股票的时候，并不限制投资哪些行业，但有的指数在挑选股票的时候，会要求只投资哪些行业的股票。

例如消费行业指数，就要求主要投资消费行业的公司，这种指数就是行业指数。而像沪深300指数，它在挑选股票的时候，并不限制行业，这种就是宽基指数。

为何要这样区分呢？因为行业指数受行业特性的影响非常大。例如，当油价大幅下跌的时候，石油开采行业的盈利就会大幅下滑，因为开采出来的石油卖不出好价格，对石油行业指数也就影响较大。但油价下跌对沪深300这种宽基指数的影响就比较小。

通常对普通投资者来说，开始投资指数基金的时候，是优先从宽基指数基金开始入手，因为这类品种相对更容易分析和理解，而行业指数基金，需要对行业有一定程度的了解，才能做好投资。

\subsubsection{市值加权和策略加权}

什么是市值加权？什么是策略加权？

市值加权比较好理解，就是股票规模越大，权重越高。

例如一个指数包含50只股票，总市值2万亿元，其中有一只市值1000亿元，另一只800亿元。那1000亿元市值的股票，在这个指数中占比就是5\%；而800亿元的股票，在这个指数中占比就是4\%。这个4\%、5\%就是权重，这就是市值加权。

而策略加权，则是按照别的方式来决定个股权重，也被称为Smart-Beta（聪明贝塔）指数。

例如红利指数，就是按照股息率来决定权重，哪只股票的股息率越高，这只股票的权重就越大。所以有的股票市值规模虽然小，但股息率高，它反而可能在红利指数中的占比更高一些。

目前从全世界范围来看，市值加权指数基金是主流，策略加权指数基金发展得晚一些，但是其品种数量和规模增长速度比市值加权指数基金更高一些。很多策略加权指数基金的长期收益也比市值加权指数基金要好。而且根据投资策略的不同，策略加权指数基金在牛熊市的表现也会有差别。

有的策略加权指数基金，在牛市涨得更多；有的策略加权指数基金，熊市跌得更少。结合不同的策略加权指数基金的特点，可以更好地帮助我们投资。

其实指数的数量是很多的，国内所有指数的数量不亚于股票的数量。一本书也不可能一一都介绍。但是在这么多的指数中，已经被开发成指数基金产品的，并不是太多。

我们掌握最核心的、最常用的指数基金品种，就足够我们做出定投计划了。所以本书这一章节，只介绍几种我们最常用到的指数，以及其对应的指数基金产品。

关于更多指数的介绍，以及其对应的指数基金的详情，可以阅读《指数基金投资指南》一书获得。

\subsection{A股常见的宽基指数：沪深300、中证500和创业板指数}
{\fangsong 沪深300指数：国内最具代表性的指数}

\subsubsection{指数简介}

沪深300指数是由中证指数公司开发的，从上海证券交易所和深圳证券交易所挑选规模最大、流动性最好的300只股票组成的。

其实指数的命名规则还挺有意思的，一般名称里的数字是多少，就代表会挑选多少只股票。比如说上证50，就包含了50只股票；沪深300，就包含了300只股票。

从市值规模上来看，沪深300指数的股票占国内股市全部规模的60\%以上，很有代表性了。在沪深300指数中，规模最小的公司的市值也在百亿元以上。

如果说哪个指数最能代表A股的综合表现，那无疑就是沪深300指数了。我们只需要看一下沪深300指数，就能知道国内股市规模最大的300只股票今天整体是涨了还是跌了。

沪深300指数的代码有两个：000300和399300。这是因为沪深300指数同时包括上海和深圳两个证券交易所的股票，所以沪深300指数在上海证券交易所的代码是000300，在深圳证券交易所的代码是399300。这两个代码其实都是代表沪深300指数的。

沪深300指数，是在2004年12月31日从1000点开始起步的。

\subsubsection{对应指数基金}

沪深300指数的影响力巨大，几乎每一家大型基金公司都有针对沪深300指数推出的指数基金，目前市面上的沪深300指数基金多达几十只，这里列举部分有代表性的沪深300指数基金品种，如表 \ref{tab:5-1} 所示（相关数据是到2019年一季度的数据，如无特殊说明，下同）。

{\scriptsize
\setlength{\tabcolsep}{2.75pt}
\renewcommand{\theHtable}{5-1}
\renewcommand{\arraystretch}{1.2}
\renewcommand{\thetable}{5.1}
\setcounter{table}{0}
\begin{longtable}{>{\raggedright\arraybackslash}m{0.28\textwidth}>{\centering\arraybackslash}m{0.13\textwidth}>{\centering\arraybackslash}m{0.12\textwidth}>{\centering\arraybackslash}m{0.15\textwidth}>{\centering\arraybackslash}m{0.12\textwidth}>{\centering\arraybackslash}m{0.12\textwidth}}
\caption{部分沪深300指数基金}\label{tab:5-1}\\[1.2ex]
\hline
\noalign{\vskip 4pt}
\multicolumn{1}{>{\centering\arraybackslash}m{0.28\textwidth}}{\rule[-1.4ex]{0pt}{5.2ex}\bfseries 基金简称} & \multicolumn{1}{>{\centering\arraybackslash}m{0.13\textwidth}}{\bfseries 基金代码} & \multicolumn{1}{>{\centering\arraybackslash}m{0.12\textwidth}}{\shortstack[c]{\bfseries 管理费率\\\bfseries （\%）}} & \multicolumn{1}{>{\centering\arraybackslash}m{0.15\textwidth}}{\shortstack[c]{\bfseries 基金规模\\\bfseries （亿元）}} & \multicolumn{1}{>{\centering\arraybackslash}m{0.12\textwidth}}{\shortstack[c]{\bfseries 成立年限\\\bfseries （年）}} & \multicolumn{1}{>{\centering\arraybackslash}m{0.12\textwidth}}{\bfseries 场内/场外} \\
\hline
\endfirsthead
\multicolumn{6}{r}{续表}\\[1ex]
\hline
\noalign{\vskip 4pt}
\multicolumn{1}{>{\centering\arraybackslash}m{0.28\textwidth}}{\rule[-1.4ex]{0pt}{5.2ex}\bfseries 基金简称} & \multicolumn{1}{>{\centering\arraybackslash}m{0.13\textwidth}}{\bfseries 基金代码} & \multicolumn{1}{>{\centering\arraybackslash}m{0.12\textwidth}}{\shortstack[c]{\bfseries 管理费率\\\bfseries （\%）}} & \multicolumn{1}{>{\centering\arraybackslash}m{0.15\textwidth}}{\shortstack[c]{\bfseries 基金规模\\\bfseries （亿元）}} & \multicolumn{1}{>{\centering\arraybackslash}m{0.12\textwidth}}{\shortstack[c]{\bfseries 成立年限\\\bfseries （年）}} & \multicolumn{1}{>{\centering\arraybackslash}m{0.12\textwidth}}{\bfseries 场内/场外} \\
\hline
\endhead
\hline
\endfoot
\hline
\multicolumn{6}{l}{\footnotesize 资料来源：Choice金融终端。}\\
\endlastfoot

易方达沪深300发起式ETF联接A & 110020 & 0.15 & 48.58 & 9.9 & 场外 \\
易方达沪深300发起式ETF & 510310 & 0.15 & 49.54 & 6.2 & 场内 \\
华夏沪深300ETF联接A & 000051 & 0.50 & 119.39 & 9.9 & 场外 \\
国寿安保沪深300ETF联接 & 000613 & 0.50 & 5.23 & 5.0 & 场外 \\
前海开源沪深300指数 & 000656 & 0.50 & 0.13 & 4.9 & 场外 \\
天弘沪深300指数A & 000961 & 0.50 & 19.22 & 4.3 & 场外 \\
广发沪深300ETF联接C & 002987 & 0.50 & 6.57 & 2.9 & 场外 \\
中金沪深300A & 003015 & 0.50 & 0.18 & 2.8 & 场外 \\
中金沪深300C & 003579 & 0.50 & 0.02 & 2.5 & 场外 \\
南方沪深300ETF联接C & 004342 & 0.50 & 0.16 & 2.2 & 场外 \\
平安300ETF联接A & 005639 & 0.50 & 3.21 & 1.1 & 场外 \\
平安300ETF联接C & 005640 & 0.50 & 4.99 & 1.1 & 场外 \\
华夏沪深300ETF联接C & 005658 & 0.50 & 9.83 & 1.3 & 场外 \\
国泰沪深300指数C & 005867 & 0.50 & 0.10 & 1.1 & 场外 \\
天弘沪深300指数C & 005918 & 0.50 & 1.76 & 1.1 & 场外 \\
华泰柏瑞沪深300ETF联接C & 006131 & 0.50 & 3.21 & 0.9 & 场外 \\
国泰沪深300指数A & 020011 & 0.50 & 20.37 & 11.5 & 场外 \\
嘉实沪深300ETF & 159919 & 0.50 & 187.43 & 7.1 & 场内 \\
南方沪深300ETF & 159925 & 0.50 & 11.50 & 6.3 & 场内 \\
嘉实沪深300ETF联接（LOF）A & 160706 & 0.50 & 168.77 & 13.7 & 场内/场外 \\
嘉实沪深300ETF联接（LOF）C & 160724 & 0.50 & 0.05 & 0.8 & 场外 \\
东吴沪深300A & 165806 & 0.50 & 0.05 & 7.2 & 场外 \\
东吴沪深300C & 165810 & 0.50 & 0.001 & 0.9 & 场外 \\
南方沪深300ETF联接A & 202015 & 0.50 & 9.92 & 10.2 & 场外 \\
广发沪深300ETF联接A & 270010 & 0.50 & 17.17 & 10.4 & 场外 \\
华泰柏瑞沪深300ETF联接A & 460300 & 0.50 & 5.05 & 7.0 & 场外 \\
工银沪深300指数A & 481009 & 0.50 & 36.27 & 10.2 & 场外 \\
汇添富沪深300指数（LOF）A & 501043 & 0.50 & 0.70 & 1.7 & 场内/场外 \\
汇添富沪深300指数（LOF）C & 501045 & 0.50 & 0.36 & 1.7 & 场内/场外 \\
华泰柏瑞沪深300ETF & 510300 & 0.50 & 332.99 & 7.1 & 场内 \\
华夏沪深300ETF & 510330 & 0.50 & 228.10 & 6.4 & 场内 \\
广发沪深300ETF & 510360 & 0.50 & 22.10 & 3.8 & 场内 \\
国寿安保沪深300ETF & 510380 & 0.50 & 5.62 & 1.3 & 场内 \\
平安沪深300ETF & 510390 & 0.50 & 42.98 & 1.4 & 场内 \\
富荣沪深300指数增强A & 004788 & 0.60 & 0.001 & 1.3 & 场外 \\
富荣沪深300指数增强C & 004789 & 0.60 & 0.13 & 1.3 & 场外 \\
农银沪深300指数C & 005152 & 0.60 & 0.0001 & 1.2 & 场外 \\
农银沪深300指数A & 660008 & 0.60 & 6.95 & 8.1 & 场外 \\
前海联合沪深300指数A & 003475 & 0.65 & 0.46 & 2.5 & 场外 \\
泰达宏利沪深300指数增强C & 003548 & 0.65 & 0.01 & 2.3 & 场外 \\
泰达宏利沪深300指数增强A & 162213 & 0.65 & 1.71 & 9.1 & 场外 \\
鹏华沪深300指数（LOF）A & 160615 & 0.75 & 3.14 & 10.1 & 场内/场外 \\
长盛沪深300指数（LOF） & 160807 & 0.75 & 0.40 & 8.8 & 场内/场外 \\
建信沪深300指数（LOF） & 165309 & 0.75 & 4.92 & 9.6 & 场内/场外 \\
大成沪深300指数A & 519300 & 0.75 & 15.98 & 13.1 & 场外 \\
创金合信沪深300指数增强A & 002310 & 0.80 & 0.73 & 3.4 & 场外 \\
创金合信沪深300指数增强C & 002315 & 0.80 & 1.96 & 3.4 & 场外 \\
安信量化精选沪深300指数增强A & 003957 & 0.80 & 0.004 & 2.2 & 场外 \\
安信量化精选沪深300指数增强C & 003958 & 0.80 & 0.74 & 2.2 & 场外 \\
易方达沪深300量化增强 & 110030 & 0.80 & 9.85 & 6.9 & 场外 \\
兴全沪深300指数（LOF）A & 163407 & 0.80 & 20.24 & 8.6 & 场内/场外 \\
西部利得沪深300指数增强A & 673100 & 0.80 & 0.17 & 2.2 & 场外 \\
国富沪深300指数增强 & 450008 & 0.85 & 1.26 & 9.7 & 场外 \\
国泰沪深300指数增强A & 000512 & 0.90 & 0.29 & 5.0 & 场外 \\
国泰沪深300指数增强C & 002063 & 0.90 & 0.20 & 3.5 & 场外 \\
博时沪深300指数C & 002385 & 0.98 & 2.77 & 3.3 & 场外 \\
博时沪深300指数A & 050002 & 0.98 & 54.08 & 15.8 & 场外 \\
长城久泰沪深300指数A & 200002 & 0.98 & 7.19 & 15.0 & 场外 \\
\end{longtable}
}



这么多指数基金可怎么选呢？

说实话，由于大家追踪的都是同一个指数，而成份股的挑选规则都一样，所以高度雷同。由此可见，指数基金具有明显的先发优势，也就是说，早期推出的沪深300指数基金，规模比后期推出的基金规模要大。

螺丝钉建议大家挑选沪深300指数基金时，遵循两个思路：第一个是寻找费用最低、误差最小的品种，这样的品种与指数的走势最接近；第二个是寻找有特色的增强型指数基金。不过，目前增强型指数基金的增强操作一般都不公开，后续能否长期保证增强收益也不确定，相比普通的指数基金还是有一定风险的。

{\fangsong 中证500指数：最容易获得超额收益的指数}

\subsubsection{指数简介}

中证500指数是由中证指数公司开发的，将沪深300指数的300家公司排除，再将最近一年日均总市值排名前300名的企业也排除，然后在剩下的公司中，选择日均总市值排名前500名的企业组成的。

从定位上来看，中证500指数以中型上市公司为主，与沪深300指数的股票没有任何重合。

中证500指数是从2004年12月31日1000点开始的，它也有两个代码：在上海证券交易所的代码是000905，在深圳证券交易所的代码是399905。

\subsubsection{对应指数基金}

跟沪深300指数一样，中证500指数的影响力也很大，也是非常出名的指数之一，所以也有很多基金公司围绕它开发指数基金。中证500指数相关的指数基金如表 \ref{tab:5-2} 所示。

{\scriptsize
\setlength{\tabcolsep}{2.75pt}
\renewcommand{\theHtable}{5-2}
\renewcommand{\arraystretch}{1.2}
\renewcommand{\thetable}{5.2}
\begin{longtable}{>{\raggedright\arraybackslash}m{0.28\textwidth}>{\centering\arraybackslash}m{0.13\textwidth}>{\centering\arraybackslash}m{0.12\textwidth}>{\centering\arraybackslash}m{0.15\textwidth}>{\centering\arraybackslash}m{0.12\textwidth}>{\centering\arraybackslash}m{0.12\textwidth}}
\caption{部分中证500指数基金}\label{tab:5-2}\\[1.2ex]
\hline
\noalign{\vskip 4pt}
\multicolumn{1}{>{\centering\arraybackslash}m{0.28\textwidth}}{\rule[-1.4ex]{0pt}{5.2ex}\bfseries 基金简称} & \multicolumn{1}{>{\centering\arraybackslash}m{0.13\textwidth}}{\bfseries 基金代码} & \multicolumn{1}{>{\centering\arraybackslash}m{0.12\textwidth}}{\shortstack[c]{\bfseries 管理费率\\\bfseries （\%）}} & \multicolumn{1}{>{\centering\arraybackslash}m{0.15\textwidth}}{\shortstack[c]{\bfseries 基金规模\\\bfseries （亿元）}} & \multicolumn{1}{>{\centering\arraybackslash}m{0.12\textwidth}}{\shortstack[c]{\bfseries 成立年限\\\bfseries （年）}} & \multicolumn{1}{>{\centering\arraybackslash}m{0.12\textwidth}}{\bfseries 场内/场外} \\
\hline
\endfirsthead
\multicolumn{6}{r}{续表}\\[1ex]
\hline
\noalign{\vskip 4pt}
\multicolumn{1}{>{\centering\arraybackslash}m{0.28\textwidth}}{\rule[-1.4ex]{0pt}{5.2ex}\bfseries 基金简称} & \multicolumn{1}{>{\centering\arraybackslash}m{0.13\textwidth}}{\bfseries 基金代码} & \multicolumn{1}{>{\centering\arraybackslash}m{0.12\textwidth}}{\shortstack[c]{\bfseries 管理费率\\\bfseries （\%）}} & \multicolumn{1}{>{\centering\arraybackslash}m{0.15\textwidth}}{\shortstack[c]{\bfseries 基金规模\\\bfseries （亿元）}} & \multicolumn{1}{>{\centering\arraybackslash}m{0.12\textwidth}}{\shortstack[c]{\bfseries 成立年限\\\bfseries （年）}} & \multicolumn{1}{>{\centering\arraybackslash}m{0.12\textwidth}}{\bfseries 场内/场外} \\
\hline
\endhead
\hline
\endfoot
\hline
\multicolumn{6}{l}{\footnotesize 资料来源：Choice金融终端。}\\
\endlastfoot

华泰柏瑞中证500ETF联接A & 001214 & 0.15 & 1.67 & 4.0 & 场外 \\
华泰柏瑞中证500ETF联接C & 006087 & 0.15 & 0.01 & 0.9 & 场外 \\
易方达中证500ETF & 510580 & 0.15 & 0.09 & 3.7 & 场内 \\
华泰柏瑞中证500ETF & 512510 & 0.15 & 3.42 & 4.0 & 场内 \\
嘉实中证500ETF联接A & 000008 & 0.15 & 10.51 & 6.2 & 场外 \\
天弘中证500指数A & 000962 & 0.50 & 7.65 & 4.3 & 场外 \\
华夏中证500ETF联接A & 001052 & 0.50 & 11.76 & 4.1 & 场外 \\
国寿安保中证500ETF联接 & 001241 & 0.50 & 1.75 & 4.0 & 场外 \\
景顺长城中证500ETF联接 & 001455 & 0.50 & 2.37 & 3.9 & 场外 \\
浙商中证500指数增强A & 002076 & 0.50 & 0.51 & 3.0 & 场外 \\
广发中证500ETF联接（LOF）C & 002903 & 0.50 & 15.23 & 2.9 & 场内/场外 \\
中金中证500A & 003016 & 0.50 & 0.25 & 2.8 & 场外 \\
中金中证500C & 003578 & 0.50 & 0.05 & 2.5 & 场外 \\
南方中证500ETF联接（LOF）C & 004348 & 0.50 & 12.81 & 2.2 & 场内/场外 \\
天弘中证500指数C & 005919 & 0.50 & 1.71 & 1.1 & 场外 \\
平安500ETF联接A & 006214 & 0.50 & 0.22 & 0.7 & 场外 \\
平安500ETF联接C & 006215 & 0.50 & 0.41 & 0.7 & 场外 \\
华夏中证500ETF联接C & 006382 & 0.50 & 0.07 & 0.7 & 场外 \\
嘉实中证500ETF联接C & 070039 & 0.15 & 0.03 & 0.7 & 场外 \\
嘉实中证500ETF & 159922 & 0.15 & 12.71 & 6.3 & 场内 \\
景顺长城中证500ETF & 159935 & 0.50 & 2.26 & 5.4 & 场内 \\
南方中证500ETF联接（LOF）A & 160119 & 0.50 & 52.65 & 9.7 & 场内/场外 \\
广发中证500ETF联接（LOF）A & 162711 & 0.50 & 10.83 & 9.5 & 场内/场外 \\
汇添富中证500指数（LOF）A & 501036 & 0.50 & 1.05 & 1.8 & 场内/场外 \\
汇添富中证500指数（LOF）C & 501037 & 0.50 & 0.65 & 1.8 & 场内/场外 \\
南方中证500ETF & 510500 & 0.50 & 334.67 & 6.3 & 场内 \\
广发中证500ETF & 510510 & 0.50 & 33.01 & 6.1 & 场内 \\
诺安中证500ETF & 510520 & 0.50 & 0.86 & 5.3 & 场内 \\
方正富邦中证500ETF & 510550 & 0.50 & 2.11 & 0.5 & 场内 \\
国寿安保中证500ETF & 510560 & 0.50 & 1.64 & 4.0 & 场内 \\
平安中证500ETF & 510590 & 0.50 & 16.61 & 1.2 & 场内 \\
华夏中证500ETF & 512500 & 0.50 & 22.06 & 4.1 & 场内 \\
富荣中证500指数A & 004790 & 0.60 & 0.001 & 1.3 & 场外 \\
富荣中证500指数C & 004791 & 0.60 & 0.11 & 1.3 & 场外 \\
鹏华中证500指数（LOF）A & 160616 & 0.75 & 2.46 & 9.3 & 场内/场外 \\
农银中证500指数 & 660011 & 0.75 & 0.74 & 7.5 & 场外 \\
创金合信中证500指数增强A & 002311 & 0.80 & 1.63 & 3.4 & 场外 \\
\end{longtable}
}



而且，中证500还是一个很特殊的指数：它是最容易获得超额收益的指数。因为目前国内中小盘股的散户比例比较高，中证500指数成份股经常出现不理性的涨跌，利用这些不理性的涨跌，可以获得超额收益。

围绕中证500指数，也出现了一系列的衍生品种。例如500低波动指数基金、500增强基金、基于中证500的量化基金等，每一类都有一些比较优秀的品种。

同样地，挑选中证500指数和沪深300指数的思路一样：第一个是寻找费用最低、规模较大、历史悠久、误差最小的品种；第二个是寻找有一定超额收益的中证500增强型指数基金。

{\fangsong 创业板指数：整体盈利比较低，不稳定}

\subsubsection{指数简介}

创业板指数是为了衡量创业板最主要的100家企业的平均表现而设立的，代码是399006。创业板指数限制了成份股的数量，只从创业板上市公司中，挑选出规模最大、流动性最好的100只股票。

创业板指数是从2010年5月31日1000点开始的。因为国内创业板历史比较短，它对应的指数历史自然也不长。从2010年到2014年，创业板处于下跌或比较平缓的状态，只有2015年上半年出现过一波牛市，不过从2016年以来进入一轮下跌周期。

创业板指数的公司整体规模较小，以中小型公司为主。这些公司的盈利大多没有进入稳定期，所以创业板的整体盈利规模比较低。公司开展新业务也更容易导致盈利大起大落。因此创业板指数更容易暴涨暴跌，投资者投资时要有心理准备。

\subsubsection{相关指数基金}

追踪创业板指数的相关指数基金如表 \ref{tab:5-3} 所示。

{\scriptsize
\setlength{\tabcolsep}{2.75pt}
\renewcommand{\thetable}{5.3}
\renewcommand{\theHtable}{5-3}
\renewcommand{\arraystretch}{1.2}
\begin{longtable}{>{\raggedright\arraybackslash}m{0.28\textwidth}>{\centering\arraybackslash}m{0.13\textwidth}>{\centering\arraybackslash}m{0.12\textwidth}>{\centering\arraybackslash}m{0.15\textwidth}>{\centering\arraybackslash}m{0.12\textwidth}>{\centering\arraybackslash}m{0.12\textwidth}}
\caption{部分创业板指数基金}\label{tab:5-3}\\[1.2ex]
\hline
\noalign{\vskip 4pt}
\multicolumn{1}{>{\centering\arraybackslash}m{0.28\textwidth}}{\rule[-1.4ex]{0pt}{5.2ex}\bfseries 基金简称} & \multicolumn{1}{>{\centering\arraybackslash}m{0.13\textwidth}}{\bfseries 基金代码} & \multicolumn{1}{>{\centering\arraybackslash}m{0.12\textwidth}}{\shortstack[c]{\bfseries 管理费率\\\bfseries （\%）}} & \multicolumn{1}{>{\centering\arraybackslash}m{0.15\textwidth}}{\shortstack[c]{\bfseries 基金规模\\\bfseries （亿元）}} & \multicolumn{1}{>{\centering\arraybackslash}m{0.12\textwidth}}{\shortstack[c]{\bfseries 成立年限\\\bfseries （年）}} & \multicolumn{1}{>{\centering\arraybackslash}m{0.12\textwidth}}{\bfseries 场内/场外} \\
\hline
\endfirsthead
\multicolumn{6}{r}{续表}\\[1ex]
\hline
\noalign{\vskip 4pt}
\multicolumn{1}{>{\centering\arraybackslash}m{0.28\textwidth}}{\rule[-1.4ex]{0pt}{5.2ex}\bfseries 基金简称} & \multicolumn{1}{>{\centering\arraybackslash}m{0.13\textwidth}}{\bfseries 基金代码} & \multicolumn{1}{>{\centering\arraybackslash}m{0.12\textwidth}}{\shortstack[c]{\bfseries 管理费率\\\bfseries （\%）}} & \multicolumn{1}{>{\centering\arraybackslash}m{0.15\textwidth}}{\shortstack[c]{\bfseries 基金规模\\\bfseries （亿元）}} & \multicolumn{1}{>{\centering\arraybackslash}m{0.12\textwidth}}{\shortstack[c]{\bfseries 成立年限\\\bfseries （年）}} & \multicolumn{1}{>{\centering\arraybackslash}m{0.12\textwidth}}{\bfseries 场内/场外} \\
\hline
\endhead
\hline
\endfoot
\hline
\endlastfoot

天弘创业板A & 001592 & 0.50 & 6.36 & 3.9 & 场外 \\
天弘创业板C & 001593 & 0.50 & 14.67 & 3.9 & 场外 \\
南方创业板ETF联接A & 002656 & 0.50 & 5.80 & 3.0 & 场外 \\
广发创业板ETF联接A & 003765 & 0.50 & 2.10 & 2.0 & 场外 \\
广发创业板ETF联接C & 003766 & 0.50 & 2.45 & 2.0 & 场外 \\
南方创业板ETF联接C & 004343 & 0.50 & 1.02 & 2.2 & 场外 \\
易方达创业板ETF联接C & 004744 & 0.50 & 9.27 & 2.0 & 场外 \\
工银创业板ETF联接A & 005390 & 0.50 & 0.33 & 1.2 & 场外 \\
工银创业板ETF联接C & 005391 & 0.50 & 1.09 & 1.2 & 场外 \\
建信创业板ETF联接A & 005873 & 0.50 & 0.23 & 0.9 & 场外 \\
建信创业板ETF联接C & 005874 & 0.50 & 0.07 & 0.9 & 场外 \\
华夏创业板ETF联接A & 006248 & 0.50 & 0.22 & 0.8 & 场外 \\
华夏创业板ETF联接C & 006249 & 0.50 & 0.14 & 0.8 & 场外 \\
博时创业板ETF联接C & 006733 & 0.50 & 0.0001 & 0.5 & 场外 \\
博时创业板ETF联接A & 050021 & 0.50 & 0.33 & 8.0 & 场外 \\
易方达创业板ETF联接A & 110026 & 0.50 & 27.90 & 7.7 & 场外 \\
博时创业板ETF & 159908 & 0.50 & 0.39 & 8.0 & 场内 \\
易方达创业板ETF & 159915 & 0.50 & 185.88 & 7.7 & 场内 \\
南方创业板ETF & 159948 & 0.50 & 7.15 & 3.0 & 场内 \\
广发创业板ETF & 159952 & 0.50 & 13.21 & 2.1 & 场内 \\
嘉实创业板ETF & 159955 & 0.50 & 0.38 & 1.9 & 场内 \\
建信创业板ETF & 159956 & 0.50 & 0.77 & 1.3 & 场内 \\
华夏创业板ETF & 159957 & 0.50 & 0.96 & 1.5 & 场内 \\
工银瑞信创业板ETF & 159958 & 0.50 & 2.31 & 1.4 & 场内 \\
国泰创业板指数（LOF） & 160223 & 0.50 & 1.13 & 2.5 & 场内/场外 \\
长城创业板指数增强发起式A & 001879 & 1.00 & 0.29 & 2.0 & 场外 \\
融通创业板指数C & 004870 & 1.00 & 0.01 & 1.9 & 场外 \\
鹏华创业板分级 & 160637 & 1.00 & 2.45 & 4.0 & 场外 \\
富国创业板指数分级 & 161022 & 1.00 & 57.38 & 5.7 & 场外 \\
融通创业板指数A & 161613 & 1.00 & 5.63 & 7.1 & 场外 \\
诺安创业板指数增强（LOF） & 163209 & 1.00 & 0.11 & 7.2 & 场内/场外 \\
\end{longtable}
}



\subsection{港股常见的宽基指数：恒生指数、H股指数和香港中小指数}
{\fangsong 恒生指数：港股最具有代表性的指数}

\subsubsection{指数简介}

国内除了上海和深圳两个证券交易所之外，还有一个非常特殊但又非常重要的证券交易所——香港证券交易所（简称港交所）。

香港在回归之前，就有几十年的证券交易历史了，比A股的历史还要长一些。香港股票市场也是一个比较成熟的股票市场。从全球股票市场的综合排名来看，港交所是全球前十的市场。港股市场是与内地关系最密切的市场之一，像我们熟悉的腾讯、比亚迪、联想等都在港交所上市交易。

恒生指数是港股最具有代表性的指数，代码是HSI，成立于1964年，历史悠久。恒生指数是由香港恒生指数公司开发的，恒生指数的点数、估值和历史估值可以在恒生指数官网上查到。

恒生指数代表的是港股的蓝筹股，由港交所所有上市公司中的50家规模最大、流动性最好的公司组成，用来反映香港股市的整体水平，单只成份股的最高比例是15\%。同时，恒生指数是一个交易指数，围绕它成立了一系列的金融衍生品，所以它也非常注重流动性。

简单来说，恒生指数就是由港交所市值规模最大、成交最活跃的50家企业组成，每季度重新选一次。这个指数和上证50指数很像，是一只蓝筹股指数。像我们熟悉的中国移动、腾讯等，都在港交所上市，它们自身的规模很大，所以也会被选入恒生指数。

恒生指数虽然说是香港股市的代表指数，但实际上现在与A股的联系越来越紧密了：恒生指数最近十几年成份股的公司中内地公司占比越来越高，香港本地公司占比逐渐下降。这也是指数的一个优势，能够自动地完成新陈代谢。特别是港股这种以机构为主的市场，定价更加成熟，企业经营不好会以更快的速度被淘汰。内地经济发展速度快，也推动了恒生指数中内地公司的占比越来越高。

\subsubsection{盈富基金}

恒生指数在中国香港本地也有对应的指数基金，例如盈富基金，就是在香港本地可以投资的一只恒生指数基金。

盈富基金很有传奇色彩。1997年亚洲金融危机爆发，当时很多亚洲国家纷纷倒在像乔治·索罗斯（George Soros）这样的金融大鳄脚下，后来索罗斯盯上了港股。

以索罗斯为首的国际金融投机家当时采用了一种立体式的投机策略。简单来说，就是做空港币，香港特别行政区政府（简称香港特区政府）为了维护汇率，被迫提高利率，利率上升后会对股市产生负面影响，再配合这种负面影响做空股票市场，将事先准备好的股票抛出，打压股指期货，做空获利。

对索罗斯来说，这场金融战胜利的关键就是使恒生指数下跌。对香港特区政府来说，自救的关键就是维护恒生指数。最后香港特区政府通过提高利率和做空成本、大量买入恒生指数成份股的方式保住了恒生指数，其中最关键的就是付出了1181亿港元买入了大量的恒生指数成份股。

香港特区政府战胜了索罗斯，但是这些股份也不能总是拿在政府手里。如果直接卖掉，可能会对当时还不稳定的香港股市再次产生冲击。所以政府就用这些股份，组织成立了一只ETF（交易型开放式指数基金），分批上市，这就是盈富基金的由来。

政府为了保证盈富基金认购完成，给盈富基金设置了非常多的有吸引力的特点，如超级低的0.03\%的管理费率，赠送红股，每年两次稳定派息等。

这些优厚的待遇，现在几乎没有哪只指数基金能复制了。如果有港股账户，那盈富基金就是恒生指数基金的绝佳选择。

不过盈富基金目前还不能通过港股通来投资，个人直接投资盈富基金比较麻烦，并且盈富基金的投资门槛比内地的基金要高很多。所以我们投资恒生指数，是以QDII型恒生指数基金为主。QDII的意思是合格境内机构投资者。我们可以把这种基金理解成一种“代购”：基金公司拿人民币，合法地投资香港及海外市场，如港股、美股、德股等。

\subsubsection{对应指数基金}

相关的指数基金如表 \ref{tab:5-4} 所示。

{\scriptsize
\setlength{\tabcolsep}{2.75pt}
\renewcommand{\theHtable}{5-4}
\renewcommand{\arraystretch}{1.2}
\renewcommand{\thetable}{5.4}
\begin{longtable}{>{\raggedright\arraybackslash}m{0.28\textwidth}>{\centering\arraybackslash}m{0.13\textwidth}>{\centering\arraybackslash}m{0.12\textwidth}>{\centering\arraybackslash}m{0.15\textwidth}>{\centering\arraybackslash}m{0.12\textwidth}>{\centering\arraybackslash}m{0.12\textwidth}}
\caption{QDII恒生指数基金}\label{tab:5-4}\\[1.2ex]
\hline
\noalign{\vskip 4pt}
\multicolumn{1}{>{\centering\arraybackslash}m{0.28\textwidth}}{\rule[-1.4ex]{0pt}{5.2ex}\bfseries 基金简称} & \multicolumn{1}{>{\centering\arraybackslash}m{0.13\textwidth}}{\bfseries 基金代码} & \multicolumn{1}{>{\centering\arraybackslash}m{0.12\textwidth}}{\shortstack[c]{\bfseries 管理费率\\\bfseries （\%）}} & \multicolumn{1}{>{\centering\arraybackslash}m{0.15\textwidth}}{\shortstack[c]{\bfseries 基金规模\\\bfseries （亿元）}} & \multicolumn{1}{>{\centering\arraybackslash}m{0.12\textwidth}}{\shortstack[c]{\bfseries 成立年限\\\bfseries （年）}} & \multicolumn{1}{>{\centering\arraybackslash}m{0.12\textwidth}}{\bfseries 场内/场外} \\
\hline
\endfirsthead
\multicolumn{6}{r}{续表}\\[1ex]
\hline
\noalign{\vskip 4pt}
\multicolumn{1}{>{\centering\arraybackslash}m{0.28\textwidth}}{\rule[-1.4ex]{0pt}{5.2ex}\bfseries 基金简称} & \multicolumn{1}{>{\centering\arraybackslash}m{0.13\textwidth}}{\bfseries 基金代码} & \multicolumn{1}{>{\centering\arraybackslash}m{0.12\textwidth}}{\shortstack[c]{\bfseries 管理费率\\\bfseries （\%）}} & \multicolumn{1}{>{\centering\arraybackslash}m{0.15\textwidth}}{\shortstack[c]{\bfseries 基金规模\\\bfseries （亿元）}} & \multicolumn{1}{>{\centering\arraybackslash}m{0.12\textwidth}}{\shortstack[c]{\bfseries 成立年限\\\bfseries （年）}} & \multicolumn{1}{>{\centering\arraybackslash}m{0.12\textwidth}}{\bfseries 场内/场外} \\
\hline
\endhead
\hline
\endfoot
\hline
\multicolumn{6}{l}{\footnotesize 资料来源：Choice金融终端。}\\
\endlastfoot

华夏沪港通恒生ETF联接（QDII）A & 000948 & 0.50 & 7.86 & 4.4 & 场外 \\
南方恒指ETF联接C & 005659 & 0.50 & 0.006 & 1.2 & 场外 \\
华夏沪港通恒生ETF联接（QDII）C & 005734 & 0.50 & 0.02 & 1.2 & 场外 \\
南方恒指ETF联接A & 501302 & 0.50 & 0.53 & 1.8 & 场内 \\
南方恒指ETF & 513600 & 0.50 & 0.69 & 4.4 & 场内 \\
华夏沪港通恒生ETF（QDII） & 513660 & 0.50 & 8.31 & 4.4 & 场内 \\
华夏恒生ETF联接（QDII）A & 000071 & 0.60 & 8.99 & 6.8 & 场外 \\
华夏恒生ETF联接现汇（QDII）A & 000075 & 0.60 & 8.99 & 6.8 & 场外 \\
华夏恒生ETF联接现钞（QDII）A & 000076 & 0.60 & 8.99 & 6.8 & 场外 \\
华夏恒生ETF联接（QDII）C & 006381 & 0.60 & 0.007 & 0.7 & 场外 \\
华夏恒生ETF（QDII） & 159920 & 0.60 & 39.00 & 6.8 & 场内 \\
大成恒生指数（QDII-LOF） & 160924 & 1.00 & 0.25 & 1.8 & 场内/场外 \\
汇添富恒生指数分级（QDII） & 164705 & 1.20 & 3.62 & 5.2 & 场内 \\
\end{longtable}
}




{\fangsong H股指数：长期被称为“价值洼地”}

\subsubsection{什么是H股指数}

H股指的是公司在内地注册，却在香港上市，用港币交易的股票。

内地企业到香港上市的很多，从1993年青岛啤酒到香港上市至今，已经有160多家企业到香港上市了。

为了衡量这些公司股票的表现，恒生指数公司编制了恒生中国企业指数，也就是我们熟知的恒生国企指数，简称H股指数。

一般人对H股指数存在一个误解，看到“国企”两个字，就觉得H股指数的成份股全是国企，但其实也是有很多民企的。不过，目前确实是以国企为主的。

最初的H股指数，只有10只成份股。从2000年开始，才改为40只成份股，并且从2000点起步，沿用至2017年。从2017年开始，H股指数进行了扩容，成份股数量增加，从原来的40只股票，变成了50只股票。

\subsubsection{为什么H股指数要扩容}

这是因为H股指数最初的规则已经无法满足投资者的需求了。

H股指数最初的目的，是反映在香港上市的内地公司的平均走势。但是很多主营业务在内地的公司，注册地不在内地，它们在香港或者在海外的其他地区注册，但也是在香港上市的。

按照H股本身的规则，这些确确实实是内地的公司就不能被纳入H股指数。这样一来，H股指数就有违它的初衷，不能够全面反映出在香港上市的内地公司的平均走势。

另外，原先的H股指数，过于集中在与金融相关的国有企业上，对民营企业纳入的比例低一些。

所以2017年，恒生指数公司决定对H股指数进行扩容。主要有以下变化。

（1）数目变多了。成份股从原来的40只股票，变成50只股票。

（2）对企业经营情况要求更高了。新的规则要求入选成份股的公司必须赢利，而且经营现金流、分红都要是正数。

（3）最主要的是增加了10只红筹股和民营企业。新增加的这10只成份股，会分批被纳入。主要是腾讯、中移动等公司，其中占比最高的就是腾讯。

从2018年3月至2019年3月，H股指数会分为5个阶段分批纳入新的成份股。

到2019年3月，扩容完成之后，这10只新加入的股票合计能占到新的H股指数的30\%左右，原H股指数的成份股占剩下的约占70\%。

这样分批纳入，可以减少对H股指数和其他成份股的冲击。

H股扩容纳入的红筹股和民营企业，包括了腾讯等优秀民营企业，相比较原来H股指数金融行业一家独大的情形，结构更加稳定，更有利于降低H股指数的长期风险。再加上扩容之后对成份股的财务等要求更加严格了，所以对H股指数长期是利好的。

H股指数扩容后，也可以更好地反映在香港上市的内地企业的实际情况。

相关的H股指数基金如表 \ref{tab:5-5} 所示。注意，H股指数基金的场外基金品种，经常与恒生指数相混淆。像有的地方就会把恒生国企指数基金简写成恒生指数基金，注意识别一下基金全称就不会混淆了。

{\scriptsize
\setlength{\tabcolsep}{2.75pt}
\renewcommand{\theHtable}{5-5}
\renewcommand{\arraystretch}{1.2}
\renewcommand{\thetable}{5.5}
\begin{longtable}{>{\raggedright\arraybackslash}m{0.28\textwidth}>{\centering\arraybackslash}m{0.13\textwidth}>{\centering\arraybackslash}m{0.12\textwidth}>{\centering\arraybackslash}m{0.15\textwidth}>{\centering\arraybackslash}m{0.12\textwidth}>{\centering\arraybackslash}m{0.12\textwidth}}
\caption{场外H股指数基金}\label{tab:5-5}\\[1.2ex]
\hline
\noalign{\vskip 4pt}
\multicolumn{1}{>{\centering\arraybackslash}m{0.28\textwidth}}{\rule[-1.4ex]{0pt}{5.2ex}\bfseries 基金简称} & \multicolumn{1}{>{\centering\arraybackslash}m{0.13\textwidth}}{\bfseries 基金代码} & \multicolumn{1}{>{\centering\arraybackslash}m{0.12\textwidth}}{\shortstack[c]{\bfseries 管理费率\\\bfseries （\%）}} & \multicolumn{1}{>{\centering\arraybackslash}m{0.15\textwidth}}{\shortstack[c]{\bfseries 基金规模\\\bfseries （亿元）}} & \multicolumn{1}{>{\centering\arraybackslash}m{0.12\textwidth}}{\shortstack[c]{\bfseries 成立年限\\\bfseries （年）}} & \multicolumn{1}{>{\centering\arraybackslash}m{0.12\textwidth}}{\bfseries 场内/场外} \\
\hline
\endfirsthead
\multicolumn{6}{r}{续表}\\[1ex]
\hline
\noalign{\vskip 4pt}
\multicolumn{1}{>{\centering\arraybackslash}m{0.28\textwidth}}{\rule[-1.4ex]{0pt}{5.2ex}\bfseries 基金简称} & \multicolumn{1}{>{\centering\arraybackslash}m{0.13\textwidth}}{\bfseries 基金代码} & \multicolumn{1}{>{\centering\arraybackslash}m{0.12\textwidth}}{\shortstack[c]{\bfseries 管理费率\\\bfseries （\%）}} & \multicolumn{1}{>{\centering\arraybackslash}m{0.15\textwidth}}{\shortstack[c]{\bfseries 基金规模\\\bfseries （亿元）}} & \multicolumn{1}{>{\centering\arraybackslash}m{0.12\textwidth}}{\shortstack[c]{\bfseries 成立年限\\\bfseries （年）}} & \multicolumn{1}{>{\centering\arraybackslash}m{0.12\textwidth}}{\bfseries 场内/场外} \\
\hline
\endhead
\hline
\endfoot
\hline
\endlastfoot

南方恒生中国企业（QDII-ETF） & 159954 & 0.50 & 1.58 & 1.3 & 场内 \\
平安港股通恒生中国企业ETF & 159960 & 0.50 & 4.64 & 0.7 & 场内 \\
建信港股通恒生中国企业ETF & 513680 & 0.50 & 2.36 & 0.4 & 场内 \\
易方达恒生中国企业ETF联接（QDII）C & 005675 & 0.60 & 0.91 & 1.3 & 场外 \\
易方达恒生中国企业ETF联接（QDII）A & 110031 & 0.60 & 12.60 & 6.8 & 场外 \\
易方达恒生中国企业ETF联接现汇（QDII）A & 110032 & 0.60 & 12.60 & 6.8 & 场外 \\
易方达恒生中国企业ETF联接现钞（QDII）A & 110033 & 0.60 & 12.60 & 6.8 & 场外 \\
易方达恒生国企ETF（QDII） & 510900 & 0.60 & 73.54 & 6.8 & 场内 \\
嘉实恒生中国企业（QDII-LOF） & 160717 & 0.75 & 2.91 & 8.7 & 场内/场外 \\
银华恒生国企指数分级（QDII） & 161831 & 1.00 & 27.82 & 5.1 & 场内 \\
\end{longtable}
}





{\fangsong 香港中小指数：港股中小盘股的代表}

\subsubsection{指数简介}

我们前面介绍过沪深300指数和中证500指数。沪深300指数代表A股最核心、规模最大的300只大盘股。中证500指数代表随后的500家最大的中盘股。

在港股市场，恒生指数和H股指数的定位，就类似于沪深300指数，投资的都是大盘股。那港股市场的中小盘股有没有对应的指数呢？答案是肯定的，那就是香港中小指数。

香港中小指数有如下特点。

（1）香港中小指数主要投资的是H股和红筹股，也就是主营业务在内地的公司。

（2）香港中小指数的中小，放在港股市场里属于中小盘股票。不过实际上，香港中小指数的平均市值规模大约是300亿元。这个规模已经是达到了沪深300指数成份股的规模了，放在A股属于大盘股范畴。

A股中盘股指数的代表是中证500指数，截至2018年年底，平均市值规模是150亿元左右。香港中小指数的平均市值规模是中证500指数的两倍左右。

（3）港股有一个比较大的风险，就是老千股。简单来说，就是港股中有一批小盘股，专门以坑普通投资者为获利方式。所以在港股投资中小盘股的时候，如果只投资个股，很容易遇到这种风险。

不过好在香港中小指数本身成份股比较分散，有150多只成份股。单只成份股的比例上限是5\%，所以即使单个公司有老千股问题，指数基金也不至于“伤筋动骨”。风险比直接投资个股要低很多。

\subsubsection{对应指数基金}

香港中小指数对应的指数基金，目前主要是一只基金，代码是501021。这也是一个LOF（上市型开放式基金），场内场外都可以投资，场内场外代码相同。

\subsection{美股常见的宽基指数：纳斯达克100指数、标普500指数}
我们介绍了中国内地的股市，介绍了中国香港股市，还有一个不容忽视的股票市场，那就是美国股市。

美国股市有两个最主要的指数：纳斯达克100指数、标普500指数。我们来分别了解一下。

{\fangsong 纳斯达克100指数}

\subsubsection{指数简介}

纳斯达克是我们比较熟悉的一个股票市场，像大家熟知的苹果（Apple）、微软（Microsoft）等公司都在纳斯达克100指数里。纳斯达克100指数投资的是纳斯达克规模最大的100家大型企业。

纳斯达克100指数的代码是NDX，从1985年100点开始。

\subsubsection{对应指数基金}

纳斯达克100指数在美国本土有很多的指数基金产品。不过我们国内的投资者投资海外市场的指数基金比较困难，主要也是通过QDII指数基金来投资纳斯达克100指数的。纳斯达克100指数相关的指数基金如表 \ref{tab:5-6} 所示。

{\scriptsize
\setlength{\tabcolsep}{2.75pt}
\renewcommand{\theHtable}{5-6}
\renewcommand{\arraystretch}{1.2}
\renewcommand{\thetable}{5.6}
\begin{longtable}{>{\raggedright\arraybackslash}m{0.28\textwidth}>{\centering\arraybackslash}m{0.13\textwidth}>{\centering\arraybackslash}m{0.12\textwidth}>{\centering\arraybackslash}m{0.15\textwidth}>{\centering\arraybackslash}m{0.12\textwidth}>{\centering\arraybackslash}m{0.12\textwidth}}
\caption{纳斯达克100指数基金}\label{tab:5-6}\\[1.2ex]
\hline
\noalign{\vskip 4pt}
\multicolumn{1}{>{\centering\arraybackslash}m{0.28\textwidth}}{\rule[-1.4ex]{0pt}{5.2ex}\bfseries 基金简称} & \multicolumn{1}{>{\centering\arraybackslash}m{0.13\textwidth}}{\bfseries 基金代码} & \multicolumn{1}{>{\centering\arraybackslash}m{0.12\textwidth}}{\shortstack[c]{\bfseries 管理费率\\\bfseries （\%）}} & \multicolumn{1}{>{\centering\arraybackslash}m{0.15\textwidth}}{\shortstack[c]{\bfseries 基金规模\\\bfseries （亿元）}} & \multicolumn{1}{>{\centering\arraybackslash}m{0.12\textwidth}}{\shortstack[c]{\bfseries 成立年限\\\bfseries （年）}} & \multicolumn{1}{>{\centering\arraybackslash}m{0.12\textwidth}}{\bfseries 场内/场外} \\
\hline
\endfirsthead
\multicolumn{6}{r}{续表}\\[1ex]
\hline
\noalign{\vskip 4pt}
\multicolumn{1}{>{\centering\arraybackslash}m{0.28\textwidth}}{\rule[-1.4ex]{0pt}{5.2ex}\bfseries 基金简称} & \multicolumn{1}{>{\centering\arraybackslash}m{0.13\textwidth}}{\bfseries 基金代码} & \multicolumn{1}{>{\centering\arraybackslash}m{0.12\textwidth}}{\shortstack[c]{\bfseries 管理费率\\\bfseries （\%）}} & \multicolumn{1}{>{\centering\arraybackslash}m{0.15\textwidth}}{\shortstack[c]{\bfseries 基金规模\\\bfseries （亿元）}} & \multicolumn{1}{>{\centering\arraybackslash}m{0.12\textwidth}}{\shortstack[c]{\bfseries 成立年限\\\bfseries （年）}} & \multicolumn{1}{>{\centering\arraybackslash}m{0.12\textwidth}}{\bfseries 场内/场外} \\
\hline
\endhead
\hline
\endfoot
\hline
\multicolumn{6}{l}{\footnotesize 资料来源：Choice金融终端。}\\
\endlastfoot

国泰纳斯达克100（QDII-ETF） & 513100 & 0.60 & 6.02 & 6.1 & 场内 \\
广发纳斯达克100指数美元现汇（QDII）A & 000055 & 0.80 & 13.84 & 4.4 & 场外 \\
大成纳斯达克100（QDII） & 000834 & 0.80 & 2.14 & 4.5 & 场外 \\
易方达纳斯达克100美元汇（QDII-LOF） & 003722 & 0.80 & 0.85 & 1.9 & 场内/场外 \\
广发纳斯达克100指数（QDII）C & 006479 & 0.80 & 0.07 & 0.6 & 场外 \\
广发纳斯达克100指数美元现汇（QDII）C & 006480 & 0.80 & 0.07 & 0.6 & 场外 \\
华安纳斯达克100指数（QDII） & 040046 & 0.80 & 1.92 & 5.8 & 场外 \\
华安纳斯达克100指数现钞（QDII） & 040047 & 0.80 & 1.92 & 5.8 & 场外 \\
华安纳斯达克100指数现汇（QDII） & 040048 & 0.80 & 1.92 & 5.8 & 场外 \\
广发纳斯达克100ETF（QDII） & 159941 & 0.80 & 0.32 & 4.0 & 场内 \\
国泰纳斯达克100指数（QDII） & 160213 & 0.80 & 7.33 & 9.1 & 场外 \\
易方达纳斯达克100人民币（QDII-LOF） & 161130 & 0.80 & 0.85 & 1.9 & 场内/场外 \\
广发纳斯达克100指数（QDII）A & 270042 & 0.80 & 13.84 & 6.8 & 场外 \\
\end{longtable}
}





{\fangsong 标普500指数}

\subsubsection{指数简介}

标普500指数历史比较悠久，也是一只蓝筹股指数，由500只成份股组成。

虽然标普500指数和沪深300指数一样是蓝筹股指数，但标普500指数并不单纯是依照上市公司的规模来选股。标普500指数并没有限制入选公司的市值规模，换句话说，标普500指数不仅有大公司（大约占90\%），还有很多中型公司（大约占10\%）。

这些公司的入选标准是其必须是一个行业的领导者。所以标普500指数是一个附带主观判断的蓝筹股指数。

除了是行业的领导者，成份股的ROE（净资产收益率）也是入选标普500指数的一个硬指标，长期ROE更好的成份股更容易入选。这也导致标普500指数的估值表现与一般市值加权的指数有一定的差异。

标普500指数也是目前追踪资金最多的一只指数，上万亿美元的基金投资在标普500指数上，远超其他美股指数。这么多的资产甚至带来了成份股买入效应：如果一只股票入选标普500指数，在入选后股票会获得大量指数基金资金的买入，导致出现一定的上涨；如果其被标普500指数淘汰，则会因为短时间卖出量比较大而下跌。

\subsubsection{标普500指数的三个特点}

第一个，更“宽”的宽基指数。

标普500指数是一个非常不错的宽基指数，与我们平时接触的沪深300、恒生等指数相比，标普500指数在行业配比上相对更加均匀。

A股沪深300指数中，金融、能源行业占比较高。这些行业周期性强，波动大，也导致沪深300指数波动比较大。

标普500指数是消费、医药、科技等弱周期性行业占比高一些，超过50\%，这使标普500指数相对来说周期性更小，走势更加平稳，更容易实现“慢牛”。

这就是标普500指数的第一个特点：是一只更加“宽”，更加“稳”的宽基指数。

第二个，代表美股市场的基准收益。

标普500指数的第二个特点是作为市场基准，反映美股市场的平均收益。换句话说，如果主动基金经理或是投资者自己投资的收益比不上标普500指数，那还不如投资标普500指数基金。巴菲特也提到，如果他去世，剩余现金会买入标普500指数基金，原因就在于标普500指数可以获取市场平均收益。

那标普500指数的收益率到底是多少呢？

标普500指数的起始点数是10，从1941年开始，不算股息收入，标普500指数在七八十年里上涨了200多倍，年化收益率大约是7\%。考虑到标普500指数历史平均股息率为2\%～2.5\%，标普500在过去的几十年里能给投资者带来9\%左右的年化收益率。

如果单看最近二三十年，美股的收益是弱于同期的港股和A股的，毕竟过去20多年是中国经济崛起的阶段，上市公司的盈利增速比美股高很多。但给人的感觉是美股牛市更多，赚到钱的人更多，原因就在于标普500指数的上升走势比较稳健，给人感觉更“牛”一些。

第三个，属于美元资产，可以分散风险。

投资标普500指数还有一个非常重要的作用，就是分散风险。一来标普500指数是美元资产，二来标普500指数自身的波动性比A股低很多，相关性也不大。所以持有一部分标普500指数基金，可以分散风险。

\subsubsection{对应指数基金}

追踪标普500指数的指数基金是目前世界上规模最大的指数基金。其中，SPY是美国本土最大的一只标普500指数基金，单只规模达到2000多亿美元，是一只庞然大物。对比一下，国内规模最大的指数基金也只有几百亿元人民币。由此可以看出，国内的指数基金市场起步还比较晚，属于比较新的事物，未来还有很大的发展空间。

我们主要是通过QDII基金来投资标普500指数，相关的指数基金如表 \ref{tab:5-7} 所示。

{\scriptsize
\setlength{\tabcolsep}{2.75pt}
\renewcommand{\theHtable}{5-7}
\renewcommand{\arraystretch}{1.2}
\renewcommand{\thetable}{5.7}
\begin{longtable}{>{\raggedright\arraybackslash}m{0.28\textwidth}>{\centering\arraybackslash}m{0.13\textwidth}>{\centering\arraybackslash}m{0.12\textwidth}>{\centering\arraybackslash}m{0.15\textwidth}>{\centering\arraybackslash}m{0.12\textwidth}>{\centering\arraybackslash}m{0.12\textwidth}}
\caption{标普500指数基金}\label{tab:5-7}\\[1.2ex]
\hline
\noalign{\vskip 4pt}
\multicolumn{1}{>{\centering\arraybackslash}m{0.28\textwidth}}{\rule[-1.4ex]{0pt}{5.2ex}\bfseries 基金简称} & \multicolumn{1}{>{\centering\arraybackslash}m{0.13\textwidth}}{\bfseries 基金代码} & \multicolumn{1}{>{\centering\arraybackslash}m{0.12\textwidth}}{\shortstack[c]{\bfseries 管理费率\\\bfseries （\%）}} & \multicolumn{1}{>{\centering\arraybackslash}m{0.15\textwidth}}{\shortstack[c]{\bfseries 基金规模\\\bfseries （亿元）}} & \multicolumn{1}{>{\centering\arraybackslash}m{0.12\textwidth}}{\shortstack[c]{\bfseries 成立年限\\\bfseries （年）}} & \multicolumn{1}{>{\centering\arraybackslash}m{0.12\textwidth}}{\bfseries 场内/场外} \\
\hline
\endfirsthead
\multicolumn{6}{r}{续表}\\[1ex]
\hline
\noalign{\vskip 4pt}
\multicolumn{1}{>{\centering\arraybackslash}m{0.28\textwidth}}{\rule[-1.4ex]{0pt}{5.2ex}\bfseries 基金简称} & \multicolumn{1}{>{\centering\arraybackslash}m{0.13\textwidth}}{\bfseries 基金代码} & \multicolumn{1}{>{\centering\arraybackslash}m{0.12\textwidth}}{\shortstack[c]{\bfseries 管理费率\\\bfseries （\%）}} & \multicolumn{1}{>{\centering\arraybackslash}m{0.15\textwidth}}{\shortstack[c]{\bfseries 基金规模\\\bfseries （亿元）}} & \multicolumn{1}{>{\centering\arraybackslash}m{0.12\textwidth}}{\shortstack[c]{\bfseries 成立年限\\\bfseries （年）}} & \multicolumn{1}{>{\centering\arraybackslash}m{0.12\textwidth}}{\bfseries 场内/场外} \\
\hline
\endhead
\hline
\endfoot
\hline
\endlastfoot

博时标普500ETF联接（QDII）A & 050025 & 0.60 & 4.14 & 6.9 & 场外 \\
博时标普500ETF（QDII） & 513500 & 0.60 & 7.02 & 5.5 & 场内 \\
易方达标普500指数美元汇（QDII-LOF） & 003718 & 0.80 & 1.44 & 2.5 & 场内/场外 \\
易方达标普500指数人民币（QDII-LOF） & 161125 & 0.80 & 1.44 & 2.5 & 场内/场外 \\
大成标普500等权重指数（QDII） & 096001 & 1.00 & 1.78 & 8.2 & 场外 \\
\end{longtable}
}




本章介绍的指数基金，有以下几个共同的特点。

（1）都是各个股票市场对应的最出名的指数，且以市值加权指数为主，也就是成份股市值规模越大，在指数中占比就越高。

（2）因为比较出名，很多指数都有多只对应的指数基金，所以可以灵活应用前面章节的内容，挑选规模合适、费率低廉、追踪误差小的品种。

不过我们也介绍过指数基金的分类。除了市值加权指数之外，还有一些非常有特色、长期收益优秀的指数基金，那就是策略加权指数基金和优秀行业指数基金。这些指数基金品种，相对要小众一些，但是通常收益会更好。掌握了这些细分品种，可以更好地提高我们定投的收益。

\subsubsection{投资者笔记}

\begin{itemize}
\item {\kaishu 指数有两种不同的分类方式，可以分为宽基指数和行业指数，也可以分为市值加权指数和策略加权指数。}
\item {\kaishu A股常见的宽基指数有：沪深300指数、中证500指数和创业板指数等。}
\item {\kaishu 港股常见的宽基指数有：恒生指数、H股指数和香港中小指数等。}
\item {\kaishu 美股常见的宽基指数有：纳斯达克100指数和标普500指数等。}
\end{itemize}
